\documentclass[10pt,a4paper]{amsart}
\usepackage[margin=19mm]{geometry}
\usepackage{amsmath,amssymb,mathtools}
\usepackage[scr=rsfs,cal=euler]{mathalfa}
\usepackage{xfrac}
\usepackage[unicode,hidelinks,bookmarks=false]{hyperref}
\newcommand{\ie}{i.e.\ }
\newcommand{\cf}{cf.\ }
\newcommand{\resp}{respectively\ }
\newcommand{\Subsection}{\subsection}
\providecommand{\nobreakdash}{\nobreak\hbox{-}}
\setlength{\emergencystretch}{2em}
\multlinegap0pt
\usepackage{bm}
\usepackage{bbm}
\usepackage{dsfont}
\usepackage{xspace}
\usepackage{cancel}
\usepackage{mathtools}
\usepackage{amsthm}
\makeatletter
\@ifundefined{theorem}{\newtheorem{theorem}{Theorem}}{}
\@ifundefined{proposition}{\newtheorem{proposition}[theorem]{\bf Proposition}}{}
\makeatother
\makeatletter
\def \G {{\Gamma}}
\def \GG {{\mathbb{G}}}
\def \HG {\mathscr{H}_{\tilde\GG}}
\newcommand{\N}{\mathbb{N}}
\def \O {{\Omega}}
\newcommand{\R}{\mathbb{R}}
\def \d {{\delta}}
\def \div {{\text{\rm div}}}
\def \e {{\varepsilon}}
\expandafter\ifx\csname giurev\endcsname\relax
\newenvironment{giurev}{\color{blue}}{\color{black}}
\fi
\def \l {{\lambda}}
\def \p {\partial}
\def \r {\varrho}
\def \rn {{\mathbb {R}}^{N}}   
\def \rnm {{\mathbb {R}}^{\tilde N}} 
\def \t {{\tau}}
\def \ti {\times}
\def \tt {\tilde}
\def \x {{\xi}}
\def \y {{\eta}}
\def \z {{\zeta}}
\makeatother
\title[Fundamental solution and Harnack inequality for subelliptic ]{Fundamental solution and Harnack inequality for subelliptic evolution operators on Carnot groups}
\author{Giulio Pecorella, Annalaura Rebucci}
\date{Source: arXiv:2509.15982v1 (CC BY 4.0)}
\begin{document}
\maketitle

\noindent\emph{Source and licence.} Fundamental solution and Harnack inequality for subelliptic evolution operators on Carnot groups. Version arXiv:2509.15982v1 (CC BY 4.0), distributed under the Creative Commons Attribution 4.0 International licence, as recorded in the arXiv licence metadata for this exact version and shown on the abstract page https://arxiv.org/abs/2509.15982v1. Full rights notice: licenses/arxiv-2509.15982-CC-BY-4.0.txt.\par\smallskip

\noindent\textit{Editorial scope.} Counted unit: the Proposition whose source label is \texttt{lifted-harnack} in the article (source0.tex lines 1744--1750, source label \texttt{lifted-harnack}), delivered together with the article's own complete original proof of it (source0.tex lines 1751--1884, one \texttt{proof} environment of 134 lines). No mathematics is written, repaired or re-derived: statement and proof are the article's own text line for line. Required context retained uncounted: Gaussian-estimates (lines 555--557); Gaussian-estimates-lower (lines 606--608); automor (lines 1461--1463); HGlifted (lines 1473--1475); tilde-meanvalue (lines 1480--1486); two-ineq (lines 1531--1537); Prop-closetoparametrix (lines 1575--1581); time-estimate (lines 1619--1621); Bound-derivative (lines 1714--1723). Every definition, display and label of those lines is retained unchanged. Omitted from this package and not redistributed: 1--554, 558--605, 609--1460, 1464--1472, 1476--1479, 1487--1530, 1538--1574, 1582--1618, 1622--1713, 1724--1743, 1885--2129. Nothing inside a retained excerpt is omitted or altered: every retained body is delivered exactly as the article's lines are written, so no omission receipt of any kind accompanies this package. Citations: the retained excerpts carry no citation command at all, so no bibliography entry of the article is transcribed and no reference is redistributed. Reference adaptation: none was needed and none was made; every reference inside the delivered unit resolves inside it, so no unresolved reference or citation is produced. Printed slips kept literal: no printed word, symbol, hypothesis or number of the counted statement or of its proof was altered. Layout and class: the article's own class is \texttt{article} with packages including framed, tikz, pstricks, amsmath, amsfonts, amssymb, mathrsfs, graphicx, xcolor, cite, soul, hyperref, mathtools, wasysym; the wrapper is the lane's pinned \texttt{amsart} editorial wrapper at 10pt on A4, which restates the theorem-like environments the retained text needs and adds the article's own macro definitions that the retained text needs (\texttt{\textbackslash G}, \texttt{\textbackslash GG}, \texttt{\textbackslash HG}, \texttt{\textbackslash N}, \texttt{\textbackslash O}, \texttt{\textbackslash R}, \texttt{\textbackslash d}, \texttt{\textbackslash div}, \texttt{\textbackslash e}, \texttt{\textbackslash l}, \texttt{\textbackslash p}, \texttt{\textbackslash r}, \texttt{\textbackslash rn}, \texttt{\textbackslash rnm}, \texttt{\textbackslash t}, \texttt{\textbackslash ti}, \texttt{\textbackslash tt}, \texttt{\textbackslash x}, \texttt{\textbackslash y}, \texttt{\textbackslash z}), with extra wrapper packages (bm, bbm, dsfont, xspace, cancel, mathtools). The delivered numbering is the unit's own. Assets: no retained span contains a figure environment, a graphics-inclusion command, a PDF-page inclusion, a table or a TikZ picture; no image file is copied into this package, the package declares no asset dependency and no image file is redistributed with it. Licence: version arXiv:2509.15982v1 of this article is distributed under the Creative Commons Attribution 4.0 International licence, as recorded in the arXiv licence metadata for this exact version and shown on the abstract page https://arxiv.org/abs/2509.15982v1. A full rights notice is delivered at licenses/arxiv-2509.15982-CC-BY-4.0.txt. Characters: every retained excerpt is pure ASCII, so the delivered engine, XeLaTeX with its default Latin Modern text font, reports no missing character.
		\begin{equation}\label{Gaussian-estimates}
			\G(x,t,\x,\t)\leq  \frac{c_T}{\big(t-\t\big)^{\frac{Q}{2}}} \;\exp{\left(-c_u\frac{d_X(x,\x)^2}{(t-\t)}\right)},
		\end{equation}

		\begin{equation}\label{Gaussian-estimates-lower}
			\G(x,t,\x,\t)\geq  \frac{\tilde c_T}{\big (t-\t\big)^{\frac{Q}{2}}} \;\exp{\left(-c_l\frac{d_X(x,\x)^2}{(t-\t)}\right)},
		\end{equation}

\begin{equation}\label{automor}
	\G_A(z,\z)= \G_A\Big(\x^{-1}\circ x,t-\t\Big)=J_A\, \G_0\Big(T_A\big(\x^{-1}\circ x	\big),t-\t\Big),
\end{equation}

\begin{equation}\label{HGlifted}
	\HG \tilde u(\tt z):=\sum_{i,j=1}^{m_1} \tilde X_i\Big(a_{ij}(z)\tilde X_j\Big)\tilde u(\tt z) + \sum_{i=1}^{m_1} b_i(z)\tilde X_i\tilde u(\tt z)+ c(z)\tilde u(\tt z)-\p_t \tilde u(\tt z),\quad \tt z\in \tilde{\O},
\end{equation}

\begin{equation} \label{tilde-meanvalue} 
	\begin{aligned}
		\tt u(\tt \z) =& \frac{1}{r} \int_{\tt\Omega_r^{(m)}(\tt \z)} \tt M_r^{(m)} (\tt \z,\tt z) \tt u(\tt z) \, d\tt z\\
		+&\frac{1}{r} \int_0^{r} \frac{\omega_m}{\r} \Biggl(
		\int_{\tt \Omega^{(m)}_\r (\tt \z)}\tt N_\r^m(\tt \z,\tt z)\, \big( \div_\GG b(z) - c(z) \big) \,\tt u(\tt z) \, d\tt z   \Biggr) \, d \r.
	\end{aligned}
\end{equation} 

	\begin{equation}\label{two-ineq}
		\begin{aligned}
			&\qquad \qquad \qquad\;\; c_\Lambda^{-1} \leq J_{A^-}, J_{A^+} \leq c_\Lambda, \\
			&c_\Lambda^{-1}  d_0(\tt x,x) \leq d_0\big(T_{A^-}(\tt x,x)\big),d_0\big(T_{A^+}(\tt x,x)\big) \leq c_\Lambda d_0(\tt x,x), \\
			& d_0\big((T_{A^-}(\tt x,x))^{-1} \,\tt\circ\, T_{A^+}(\tt x,x)\big) \leq c_\Lambda \|A^+ - A^-\|^{\frac{1}{\kappa}} d_0(\tt x,x),
		\end{aligned}
	\end{equation}

\begin{proposition}\label{Prop-closetoparametrix}
	Assume that \textbf{(H1)}, \textbf{(H2)}, \textbf{(H3)} and \textbf{(H4)} hold. Then, for every $\y>0$ there exists $K>0$ (depending only on $\y$ and structural assumptions) such that 
	\begin{equation}\label{close-to-parametrix}
		(1-\y)\tt Z(\tt z,\tt \z)\leq \tt\G(\tt z,\tt \z)\leq (1+\y)\tt Z(\tt z,\tt \z),
	\end{equation}
	for every $\tt z,\tt \z\in O_K:=\Big\{\tt Z(\tt z,\tt \z)\geq K\Big\}$.
\end{proposition}

	\begin{equation}\label{time-estimate}
		O_K\subset \Big\{(\tt x,x, t)\in \rnm\ti \rn\ti\R : 0\leq t\leq T\Big\}.
	\end{equation}

\begin{proposition}\label{Bound-derivative}	
	Assume that \textbf{(H1)},\textbf{(H2)},\textbf{(H3)} and \textbf{(H4)} hold. Then, there exist $c,K>0$ (depending only on structural assumptions) such that 
	\begin{equation}\label{bound-derivative}
		\begin{aligned}
			&|\tt X_i \tt \G(\tt z,\tt \z)|\leq c \left(\frac{1}{\sqrt{t-\t}}d_{\tt\GG}(\tt \x,\x,\tt x,x) + 1 \right) \tt \G(\tt z,\tt \z),\\
			&|\tt X_i \tt \G(\tt \z,\tt z)|\leq c \left(\frac{1}{\sqrt{t-\t}}d_{\tt\GG}(\tt x,x,\tt \x, \x) + 1 \right) \tt \G(\tt z,\tt \z),
		\end{aligned}
	\end{equation}
	for every $\tt z,\tt \z \in \tt O_K:=\big\{\tt\G(\tt z,\tt \z)\geq K\big\}$. 
\end{proposition}

\begin{proposition}\label{lifted-harnack}
	Let $\HG$ be a differential operator of the form \eqref{HGlifted} satisfying assumptions \textbf{(H1)}, \textbf{(H2)}, \textbf{(H3)} and \textbf{(H4)}. For every $m\in\N$ with $m>2$, there exists three positive constants $r_0,\e_0$ and $C_P$ (only depending on structural assumptions and $m$) such that the following holds: for every $(\tt x_0,x_0, t_0)\in\tt\O$, $0<r\leq r_0$, $0<\e\leq \e_0$ and $\tt\O^{(m)}_{5r}(\tt x_0, x_0, t_0)\subset \tt\O$ we have
	\begin{equation}\label{parabolic-Harnack-lifted}
		\sup_{\tt K_{r,\e}^{(m)}(\tt x_0, x_0, t_0)} \tt u \leq C_P \, \tt u(\tt x_0, x_0, t_0),
	\end{equation}
	for every $\tt u\geq0$ classical solution to $\HG \tt u =0$ in $\tt\O$.
\end{proposition}

\begin{proof}
	Let $m >2$, $(\tt x_0,x_0, t_0)\in\tt\O$ and $r >0$ such that $\tt \O^{(m)}_{4 r}\big(\tt x_0,x_0, t_0\big) \subset \tt \O$. We show that there exist $m^-,m^+>0$ and $r_0,\e_0,\vartheta\in (0,1)$ such that, for every $r\leq r_0$ and $\e\leq \e_0$, the following properties hold
	\begin{enumerate}
		\item $\tt K_{r,\e}^{(m)}(\tt x_0, x_0, t_0)\neq \emptyset$; \label{item1H}
		\item $ \tt M_{\vartheta r}^{(m)} \big(\tt \x,\x,\t,\tt x,x,t\big) \leq m^+r^{\frac{m-2}{m+\tt Q}}$ for every $(\tt x, x,t)\in \tt \O^{(m)}_{\vartheta r}\big(\tt \x, \x, \t\big)$; \label{item2H}
		\item $\tt \O^{(m)}_{\vartheta r}\big(\tt \x,\x, \t\big)\subset \tt \O^{(m)}_{4 r}(\tt x_0, x_0, t_0)\cap \Big\{\t\leq t_0 - \e\,r^2\Big\}$ for every $\big(\tt \x, \x,\t\big)\in K_{r,\e}^{(m)}(\tt x_0, x_0, t_0)$;\label{item3H}
		\item $\tt M_{5r}^{(m)} \big(\tt x_0, x_0,t_0,\tt x, x,t\big) \geq m^- r^m$ for every $\big(\tt x, x,t\big)\in \tt \O^{(m)}_{4r}(\tt x_0, x_0, t_0)\cap \Big\{t\leq t_0 - \e\,r^2\Big\}$.\label{item4H}
	\end{enumerate}
	Once these properties have been established, the proof of \eqref{parabolic-Harnack-lifted} is straightforward. Indeed, if $\div_{\GG} b-c=0$, for every $\big(\tt\x,\x,\t\big)\in\tt K_r^{(m)}(\tt x_0, x_0, t_0)$  we get
	\begin{equation}\label{simple-harnack}
		\begin{aligned}
			\tt u\big(\tt\x,\tt \x,\t\big)\overset{\eqref{tilde-meanvalue}}=&\frac{1}{\vartheta r } \int_{\tt\Omega_{\vartheta r}^{(m)}(\tt\x,\x,\t)} \tt M_{\vartheta r}^{(m)} \big(\tt\x,\x,\t,\tt x, x,t\big) \tt u(\tt x,x,t) \, d\tt x d x dt \\
			\overset{(\ref{item2H}.)}\leq\;& \frac{m^+  r^{\frac{m-2}{m+\tt Q}} }{\vartheta r} \int_{\tt\Omega_{\vartheta r}^{(m)}(\tt \x,\x,\t)} \tt u(\tt x, x,t) \, d\tt x d x dt\\
			\overset{(\ref{item3H}.)}\leq\;& \frac{m^+r^{\frac{m-2}{m+\tt Q}} }{\vartheta r} \int_{\tt \O^{(m)}_{4 r}(\tt x_0, x_0, t_0)\cap \left\{\t\leq t_0 - \e\,r^2\right\}} \tt u(\tt x, x,t) \, d\tt x dx dt\\
			\overset{(\ref{item4H}.)}\leq\;& \frac{5m^+ r^{\frac{m-2}{m+\tt Q}} }{\vartheta m^- r^m }\frac{1}{5 r } \int_{\tt \O^{(m)}_{5r }(\tt x_0, x_0, t_0)} \tt M_{5r}^{(m)}  \big(\tt x_0, x_0,t_0,\tt x, x,t\big) \tt u(\tt x, x,t) \, d\tt x d x dt\\\overset{\eqref{tilde-meanvalue}}\leq&\frac{5m^+  r_0^{\frac{m-2}{m+\tt Q}-m} }{\vartheta m^-} \tt u(\tt x_0, x_0,t_0),
		\end{aligned}
	\end{equation}
	and \eqref{parabolic-Harnack-lifted} follows with $C_P = \frac{5m^+  r_0^{\frac{m-2}{m+\tt Q}-m} }{\vartheta m^-}$.
	
	\medskip
	
	We now remove the assumption $\div_{\GG} b-c=0$. From \textbf{(H4)} we infer that there exists $k>0$ such that $|\div_{\GG} b-c|\leq k:=(m_1+1)M_1$. We introduce the auxiliary functions
	\begin{equation}
		\overline u(\tt x,x,t):=\mathrm{e}^{k(t-t_0)}\tt u(\tt x, x,t),\qquad \underline u(\tt x, x,t):=\mathrm{e}^{-k(t-t_0)}\tt u(\tt x,x,t).
	\end{equation}
	Since $\HG \tt u = 0$, we have that
	\begin{equation}
		\overline{\HG} \overline u: = \HG \overline u + k \overline u = 0,\qquad \underline{\HG} \underline u: = \HG \underline u - k \underline u = 0.
	\end{equation}
	We observe that $\overline{\HG}$ and $\underline{\HG}$ are operators of the form \eqref{HGlifted} with zero-order term $\overline{c}:=c+k$ and $\underline{c}:=c-k$, respectively. It is clear that \emph{\ref{item1H}.}-\emph{\ref{item4H}.} hold with the same constants $r_0,\e_0,\vartheta, m^-,m^+$. Moreover, by Gaussian estimates we infer that there exists $T,r^*>0$ such that $0<\t-t_0<T$ for every $\big(\x,\tt \x,\t\big)\in\tt\Omega_{5 r}^{(m)}(x_0,\tt x_0,t_0)$, provided $r<r^*$. Assuming without loss of generality that $r_0 \leq  r^*$, we infer that there exists $\overline{c}, \underline{c}>0$ such that 
	\begin{equation}\label{soprasottobarrato}
		\overline{c}	\tt u\big(\tt \x, \x,\t\big) \leq \overline{u}\big(\tt\x,\x,\t\big) \leq \,\tt u\big(\tt \x,\x,\t\big),\qquad \,\tt u\big(\tt \x, \x,\t\big) \leq \underline{u}\big(\tt \x, \x,\t\big) \leq \,\underline{c}\tt u\big(\tt \x, \x,\t\big),
	\end{equation} 
	for every $\big(\x,\tt \x, \t\big)\in \tt\Omega_{5 r}^{(m)}(x_0,\tt x_0,t_0)$. Let $\overline M_{r}^{(m)}$ and $\underline M_{r}^{(m)}$ denote the mean value formula kernels relevant to $\overline{\HG}$ and $\underline{\HG}$, respectively, and by $\underline\Omega_{r}^{(m)}$ and $\overline\Omega_{r}^{(m)}$ the corresponding super-level sets. Since $\div_{\GG} b - \overline{c}=\div_{\GG} b -c-k \leq 0$, the first two lines in \eqref{simple-harnack} become
	\begin{equation}
		\begin{aligned}
			\tt u\big(\tt \x, \x,\t\big) &\overset{\eqref{soprasottobarrato}}\leq\frac{1}{\overline{c}}\overline{u} \big(\tt \x,\x, \t\big)\\ 
			&\;\overset{\eqref{tilde-meanvalue}}\leq \frac{1}{\vartheta r \overline{c}} \int_{\overline{\Omega}_{\vartheta r}^{(m)}(\tt\x,\x,\t)} \overline{M}_{\vartheta r}^{(m)} \big(\tt\x,\x,\t,\tt x, x,t\big) \overline{u}(\tt x,x,t) \, d\tt x d x dt \\		
			&\;\;\overset{(\ref{item2H}.)}\leq\frac{m^+ r^{\frac{m-2}{m+\tt Q}} }{\vartheta r \overline{c}} \int_{\overline\Omega_{\vartheta r}^{(m)}(\tt\x,\x,\t)}  \overline u(\tt x, x,t) \, d\tt x  d x dt =: I.
		\end{aligned}
	\end{equation}
	Since \emph{\ref{item3H}.} holds also in the form $\overline\Omega_{\vartheta r}^{(m)}(\tt \x,\x, \t)\subset \underline \O^{(m)}_{4 r}(\tt x_0, x_0, t_0)\cap \big\{\t\leq t_0 \!-\!\e\, r^2\big\}$,  by using $\div_{\GG} b - \underline{c} =\div_{\GG} b-c+k\geq 0$ we conclude
	\begin{equation}
		I\leq \frac{5m^+r^{\frac{m-2}{m+\tt Q}}}{\vartheta m^-r^m }\,\frac{\underline{c}}{\overline{c}}\,\frac{1}{5 r } \int_{\underline \O^{(m)}_{5r }(\tt x_0, x_0, t_0)} \!\!\!\!\!\!\!\!\!\!\!\!\!\!\!\!\!\!\!\!\underline M_{5r}^{(m)}\big(\tt x_0, x_0,t_0,\tt x, x,t\big) \underline u(\tt x, x,t) \, d\tt x d x dt \leq \frac{5m^+}{\vartheta m^-} \,\frac{\underline{c}}{\overline{c}}\, r_0^{\frac{m-2}{m+\tt Q}-m} \,\tt u(\tt x_0, x_0,t_0),
	\end{equation}
	where the last inequality follows again from \eqref{tilde-meanvalue} and \eqref{soprasottobarrato}.  Hence, inequality \eqref{parabolic-Harnack-lifted} follows with $C_P = \frac{5m^+}{\vartheta m^-} \frac{\underline{c}}{\overline{c}}\, r_0^{\frac{m-2}{m+\tt Q}-m}$. 
	
	\medskip
	
	To conclude the proof of Proposition \ref{lifted-harnack}, we are now left with proving claims \emph{\ref{item1H}.}-\emph{\ref{item4H}.}
	
	
	\medskip
	\emph{\ref{item1H}.} We prove that the set 
	\begin{equation}
		\tt K_{r,\e}^{(m)}(\tt x_0, x_0, t_0)=\overline{\biggl\{\frac{\tt \Gamma(\tt x_0, x_0, t_0,\tt x, x,t)}{\big( 4\pi(t_0-t)\big)^{m/2}}  >\frac{1}{r} \biggr\}}\cap \Big\{t\leq t_0 - \e\, r^2\Big\},
	\end{equation}
 	is non-empty for every $r,\e\leq 1$. From the Gaussian lower bound \eqref{Gaussian-estimates-lower} we infer that there exists $c,c_l>0$ (depending on structural assumptions) such that $\tt K_{r,\e}^{(m)}(\tt x_0, x_0, t_0)\subset \mathcal K$, where
	\begin{equation}
		\mathcal K:=\overline{\Biggl\{\frac{c}{\big (t_0-t\big)^{\frac{\tt Q}{2}}} \;\exp{\left(-c_l\frac{d_{\tt\GG}(\tt x_0, x_0,\tt x, x)^2}{(t_0-t)}\right)}\frac{1}{\big( 4\pi(t_0-t)\big)^{m/2}}  >\frac{1}{r} \Biggr\}}\cap \Big\{t\leq t_0 - \e\,r^2\Big\}. 
	\end{equation}
	We now show that $\mathcal K$ is non-empty. Let $t:=t_0-\e\,r^2$, and notice that any $(\tt x, x)\in \rnm\ti\rn$ satisfying
	\begin{equation}
		d_{\tt\GG}(\tt x_0, x_0,\tt x, x)\leq \frac{1}{\e\,r^2c_\l}\ln\left(\frac{(4\pi)^\frac{m}{2}}{c}\big(\e\,^2r\big)^{m+\tt Q-1}\right),
	\end{equation}
	belongs to $\mathcal K$. Therefore $\mathcal K$ is non-empty, and so is $\tt K_{r,\e}^{(m)}(\tt x_0, x_0, t_0)$.
	
	
	\medskip
	\emph{\ref{item2H}.} We show that there exist $r_0,m^+>0$ such that, for every $r\leq r_0$, it holds 
	\begin{equation}
		\tt M_{ r}^{(m)}\big(\tt \x,\x,\t,\tt x,x,t\big)\leq m^+ r^{\frac{m-2}{m+\tt Q}},\qquad \text{for every }\, (\tt x, x,t)\in \tt \O^{(m)}_{r}\big(\tt \x, \x, \t\big).
	\end{equation}
	Recall that the explicit formula for the kernel $\tt M_{ r}^{(m)}$ is given by 
	\begin{equation}\label{kernel-formula}
		\begin{aligned}
			\tt M_{ r}^{(m)}\big(\tt \x,\x,\t,\tt x, x,t\big)&=\frac{2^{m} m\omega_m}{m+2}(\t-t)^\frac{m-2}{2}\left(\ln\left( r\, \frac{\tt \Gamma\big(\tt \x, \x, \t,\tt x, x,t\big)}{\big( 4\pi(\t-t)\big)^{m/2}} \right)\right)^\frac{m+2}{2}\\
			&\quad+2^m \omega_m \left((\t-t)\ln\left( r\, \frac{\tt \Gamma\big(\tt \x, \x, \t,\tt x, x,t\big)}{\big( 4\pi(\t-t)\big)^{m/2}} \right)\right)^\frac{m}{2}\cdot\\
			&\qquad\qquad \qquad \cdot\frac{\langle A(x,t)\nabla_{\tt \GG} \tt \Gamma\big(\tt \x, \x, \t,\tt x, x,t\big),\nabla_{\tt \GG}\tt \Gamma\big(\tt \x, \x, \t,\tt x, x,t\big) \rangle}{{\tt \Gamma\big(\tt \x, \x, \t,\tt x, x,t\big)}^2}\\
			&=:K_1+K_2.
		\end{aligned}
	\end{equation}
	We first focus on the term $K_2$. To this end, we show that there exist constants $r_0>0$ and $c>0$ such that, for every $r\leq r_0$, the following inequality holds 
	\begin{equation}\label{bound-kernel}
		\frac{\langle A(x,t)\nabla_{\tt \GG} \tt \Gamma\big(\tt \x, \x, \t,\tt x, x,t\big),\nabla_{\tt \GG}\tt \Gamma\big(\tt \x,\x, \t,\tt x, x,t\big) \rangle}{{\tt \Gamma\big(\tt \x, \x, \t,\tt x, x,t\big)}^2}\leq \frac{c}{(\t-t)},\qquad (\tt x, x,t)\in \tt \O^{(m)}_{r}\big(\tt \x, \x, \t\big).
	\end{equation}
	To prove \eqref{bound-kernel} we apply Proposition \ref{Bound-derivative}.  We claim that this proposition remains valid when replacing  $\big\{\tt\G\geq K\big\}$  with the level set $\tt\Omega_{1/K}^{(m)}$. Indeed, Proposition \ref{Bound-derivative} applied to  $\tt \G^{(m)}$ on the set $\Big\{\tt\G^{(m)}(0,\tt \x,\x,\t,\tt x,x,y,t\Big)\geq K\big\}\supseteq\tt\Omega_{1/K}^{(m)}\big(\tt \x,\x,\t\big)$ implies that
	\begin{equation}
		\begin{aligned}
			H^{(m)}(0,\t,0,t)\tt X_i\tt\G\big(\tt \x, \x, \t,\tt x, x,t\big)& = X_i\tt\G^{(m)}(0,\tt \x, \x, \t,0,\tt x, x,y,t)\\
			&\leq  c \left(\frac{1}{\sqrt{t-\t}}d_{\tt\GG}\big(\tt \x,  \x,\tt x,x \big) + 1 \right) \tt \G^{(m)}\big(0,\tt \x,\x ,\t,0,\tt x, x ,t\big)\\
			&= c \left(\frac{1}{\sqrt{t-\t}}d_{\tt\GG}\big(\tt \x,  \x,\tt x, x\big) + 1 \right) H^{(m)}(0,\t,0,t)\tt\G\big(\tt \x, \x, \t,\tt x, x,t\big),				
		\end{aligned}
	\end{equation}
	proving the claim. Let $K$ be the constant given by the aforementioned proposition, and set $r_0:=1/K$. By the ellipticity assumption given by \textbf{(H4)} and Proposition \ref{Bound-derivative}, we finally conclude the validity of estimate \eqref{bound-kernel}.
	
	
	We focus now on $K_1$. By Proposition \ref{Prop-closetoparametrix} and \eqref{time-estimate}, we deduce that, if we pick $r_0$ small enough, we can choose $T>0$ such that
	\begin{equation}
		\tt\Omega_{r}^{(m)}(\tt\x,\x,\t)\subseteq \Big\{(\tt x, x, t)\in \rnm\ti\rn\ti\R : 0\leq \t-t\leq T\Big\}, \qquad \text{for every }\,r\leq r_0.
	\end{equation}
	Then, by \eqref{Gaussian-estimates} we infer that there exists $c=c(r_0)>0$ such that
	\begin{equation}\label{Pozzallo}
		K_1\leq \frac{2^m m \omega_m}{m+2}(\t-t)^\frac{m}{2}\left(\ln\Bigg(\frac{c\,r}{(\t-t)^\frac{m+\tt Q}{2}}\Bigg) \right)^\frac{m+2}{2}.
	\end{equation}
	Combining \eqref{kernel-formula}, \eqref{bound-kernel}, and \eqref{Pozzallo}, the validity of \emph{\ref{item2H}.} follows.
	
	
	\medskip
	\emph{\ref{item3H}.} The key ingredient to prove this inclusion is the parabolic dilation relevant to $\tt\GG$. Indeed, by Proposition \ref{Prop-closetoparametrix} we infer that there exist $r_0,\e_0>0$ such that for every $r\leq r_0$ and $\e\leq \e_0$ it holds
	\begin{equation}\label{casinalbo}
		\tt K_{r,\e}^{(m)}(\tt x_0, x_0, t_0)\subset \overline {\tt \O^{(m)^*}_{3 r}(\tt x_0, x_0, t_0)}\cap \Big\{t\leq t_0 -\e\, r^2\Big\},
	\end{equation} 
	where
	\begin{equation}
		\tt \O^{(m)^*}_{r}(\tt x_0,x_0, t_0):=\Bigg\{(\tt x, x ,t) \in \rnm\ti\rn\ti \R :\frac{\tt Z(\tt x_0, x_0, t_0,\tt x, x ,t)}{\big( 4\pi(t_0-t)\big)^{m/2}}>\frac{2}{r} \Bigg\},
	\end{equation}
	with $\tt Z$ being the parametrix of $\tt \G$. We claim that there exists $\vartheta \in (0,1)$ such that 
	\begin{equation}\label{parabolic-inclusion}
		\tt \O^{(m)^*}_{3\vartheta r}(\tt x, x, t)\subset \tt \O^{(m)^*}_{4 r}(\tt x_0, x_0, t_0)\cap \Big\{\t\leq t_0 - \e\,r^2\Big\},
	\end{equation}
	for every $(\tt x,x, t)\in  \overline {\tt \O^{(m)^*}_{3 r}(\tt x_0, x_0, t_0)}\cap \left\{t\leq t_0 - \e\, r^2\right\}$. The validity of \emph{\ref{item3H}.} follows then by Proposition \ref{Prop-closetoparametrix}, that implies $\tt \O^{(m)^*}_{4r}(\tt x_0,x_0, t_0)\subset\tt \O^{(m)}_{4r}(\tt x_0, x_0, t_0)$. To prove claim \eqref{parabolic-inclusion}, we employ the parabolic dilation relevant to $\tt\GG=\Big(\rnm\times \rn,\tt \circ,\tt \d_r\Big)$: by the homogeneity of the parametrix we infer that for every positive $k$ it holds 
	\begin{equation}\label{scaling}
		\Big((\tt x_0,x_0)\;\tt\circ \;\tt{\d_{r}}(\tt x,x),t_0+r^2t\Big)\in \tt \O^{(m)^*}_{k r}(\tt x_0,x_0, t_0)\iff \Big((\tt x_0,x_0)\;\tt\circ \;(\tt x,\tt x),t_0+t\Big)\in \tt \O^{(m)^*}_{k}(\tt x_0, x_0, t_0).
	\end{equation}
	Hence, it is sufficient to show \eqref{parabolic-inclusion} whenever $r=1$. Let $d(\tt x_0, x_0,t_0)$ denote the distance between the compact set $\overline {\tt \O^{(m)^*}_{3}(\tt x_0, x_0, t_0)}\cap \left\{t\leq t_0 - \e\right\}$ and the boundary of $\tt \O^{(m)^*}_{4}(\tt x_0, x_0, t_0)$. Clearly, this defines a positive function continuously depending on $(\tt x_0, x_0,t_0)$. In particular, it depends on the  automorphism $T_{A(\tt x_0, x_0,t_0)}$, associated with the matrix $A$ evaluated at the point $(\tt x_0, x_0,t_0)$ (see \eqref{automor}). By \eqref{two-ineq}, we infer that there exists $d_0,\e_0>0$ (depending only on structural assumption) such that $d(\tt x_0, x_0,t_0)\geq d_0$ for every $(\tt x_0, x_0,t_0)\in \tt\O$ and $\e\leq\e_0$. Since the diameter of $\tt\O_1^{(m)^*}(\tt x, x,t)$ is uniformly bounded, by \eqref{scaling} we deduce that it is possible to find $\vartheta \in (0,1)$ such that the diameter of $\tt\O_\vartheta^{(m)^*}(\tt x,x,t)$ is not greater than $d_0$, and then we conclude the proof of the claim.
	
	
	\medskip
	\emph{\ref{item4H}.} We recall again that the explicit expression of the kernel is given by \eqref{kernel-formula}. By the ellipticity assumption \textbf{(H4)} we infer $K_2\geq0$. To prove the estimate is therefore sufficient to show that $K_1\geq m^-r^m$. This follows once we set  
	\begin{equation}
		m^-:=\frac{2^mm\omega_m}{4(m+2)}\left(\ln(5/4) \right)^\frac{m+2}{2}.
	\end{equation}
\end{proof}

\end{document}
